\section{Finite formulas for the two flows}\label{sec:numerics}

We now turn the probability transport and the stationary centering flow
into finite formulas in conditional means.  The formulas must meet two
requirements at once.  Their error on the prescribed input law must be
small, and each new mean must be evaluated at a center that is an affine
expression in earlier mean values, with scalar coefficients and controlled
size.  The latter property permits the recursive implementation in the
next section.

The main conclusion is Theorem~\ref{thm:numerical-interface}.  Its proof
has three numerical steps.  Global Picard iteration discretizes
probability transport without making the dependency depth proportional
to the number of time cells.  Short harmonic trajectories then give a
formula for the vector $K(y)^TG$ in the centering velocity.  Finally,
Picard iteration discretizes that stationary flow.  Two cancellations
are needed in the middle step: a zero derivative removes the singularity
at the lower endpoint of the kernel integral, and retaining only a
forward transport \emph{correction} cancels its apparent inverse power
of the primitive scale $r$.

Polynomial collocation of Picard iteration was developed for sampling
by \citet[Algorithm~2 and Theorem~2.3]{lee2018odes}; Gaussian smoothing
and high-order Picard approximation also appear in
\citet[Appendix~E, Algorithm~E.1 and Theorem~E.2]{chen2026picard}.
Here we need, in addition, scalar affine formulas whose dependency depth
can be fixed before their interpolation order.  An exact conditional
mean is an analytical node throughout this section.  The next section
replaces these nodes by finite oracle computations.

Fix an integer $K_0\ge2$ and $t>0$ such that
\begin{equation}\label{eq:numerical-parameters}
 40(K_0+2)t<\frac14,\qquad
 R=\delta^t,\qquad \psi=\delta^{4t},\qquad h_0=\delta^w,\qquad 0<w<t.
\end{equation}
The iteration count will first fix a depth bound.  We then choose $w$
and finally the interpolation orders; increasing those orders will not
change the depth bound.  For the centering flow set $s=\sigma/2$, where
$\sigma>0$ is the parent mean-noise parameter.  Thus the centering
formula itself will approximate noise of standard deviation $\sigma/2$.
Throughout its numerical analysis assume
\begin{equation}\label{eq:numerical-size}
 \max\{l,L/\sigma\}\le\polylog(d)\delta^b,\qquad b\ge\frac12.
\end{equation}
For probability transport alone only the corresponding bound on $l$ is
required.  All probability-time intervals lie in
$[0,(1-cR^2)^{1/2}]$, for a fixed $c>0$.  Constants may depend on fixed
interpolation orders and fixed parameters, but not on the center,
dimension, primitive scale $r$, noise parameter, or oracle.  A symbol
$B_d$ below denotes a bound $C(1+\log d)^C$ with precisely this uniformity;
the fixed exponent and constant in a size hypothesis are part of the
fixed parameters.  All assertions are for sufficiently large $d$.

\subsection{Probability transport by global Picard iteration}

Partition a probability-time interval using equal steps of size at most
$h_0$ in $-\log(1-\rho)$.  For a nonzero interval take the ceiling of
its transformed length divided by $h_0$ as the number of cells.
Cell lengths $\ell_c$ then obey
$\ell_c\le C h_0(1-\rho)$ on that cell, and the number of cells is
\[
 O\bigl(1+h_0^{-1}(1+|\log R|)\bigr).
\]
A zero-length interval requires no operations.  On each cell interpolate
at $m\ge2$ distinct equally spaced nodes, including both ends, and
integrate the resulting degree-$(m-1)$ polynomial.  Include partial
integrals to every state node.  By scaling the fixed Lagrange basis on
$[0,1]$, its absolute value sum is at most $C_m$, and its order-$k$
derivative absolute sum is at most $C_m\ell_c^{-k}$.
Consequently every partial-integral row has absolute weight sum at most
$C_m$, uniformly in the number of cells.

Let $w_{ij}$ be these partial-integral weights from the starting time
$\rho_{\mathrm{in}}$ to node $\rho_i$, with signed weights for backward
integration.  Duplicate shared endpoints as labeled nodes if convenient.
Given a starting position $y_{\mathrm{in}}$, use simultaneous updates
\begin{equation}\label{eq:probability-update}
 y_i^{[0]}=y_{\mathrm{in}},\qquad
 y_i^{[k+1]}=y_{\mathrm{in}}-
 r\sum_j w_{ij}M_{\rho_j}(y_j^{[k]}).
\end{equation}
This iteration integrates over the whole interval at every layer; it does
not advance successively from cell to cell.

\begin{proposition}[Probability-flow circuit]
\label{prop:probability-circuit}
If the input position has marginal law $\xi_{\rho_{\mathrm{in}}}$, then
after $N_p$ iterations the maximum, over all state nodes, of the RMS
error relative to the exact transport is at most
\begin{equation}\label{eq:probability-picard}
 C_m\polylog(d)\,l\sqrt d\,h_0^m
       +(C_ml)^{N_p}CrD.
\end{equation}
Every computed state is globally Lipschitz in its initial position with
constant $C_m$.  After subtracting the initial position, its Lipschitz
constant is at most $C_ml$.
\end{proposition}

\begin{proof}
Write $Y_\rho$ for the exact path from the given random input and
$Q_\rho=M_\rho(Y_\rho)$.  The marginal identity in
Proposition~\ref{prop:transport} and \eqref{eq:path-mean} give
\[
 \|Q_\rho^{(m)}\|_2\le B_dL\sqrt d(1-\rho^2)^{-m}.
\]
On a cell $I=[a_c,b_c]$ of length $\ell_c$, let $P$ be the
Taylor polynomial of $Q$ of degree $m-1$ at $a_c$.  Minkowski's
inequality applied to the integral remainder gives, at each $v\in I$,
\begin{equation}\label{eq:numerical-taylor-remainder}
 \|Q_v-P(v)\|_2
 \le \frac1{(m-1)!}\int_{a_c}^v
       (v-z)^{m-1}\|Q_z^{(m)}\|_2\,dz
 \le C_mB_dL\sqrt d\,h_0^m.
\end{equation}
The last step uses $\ell_c\le Ch_0(1-z)$ on that cell and
$1-z^2\ge1-z$.  The same estimate holds at the interpolation nodes.
Because interpolation reproduces $P$ and its basis has bounded absolute
sum, $\|Q_v-\mathcal I_cQ(v)\|_2$ obeys the same bound.
Integrate this pointwise estimate over the cells, whose total length
is at most one.  This proves the required quadrature error, for partial
integrals and for either orientation.  No supremum of the random path
has been taken.

Let $e_k=\max_i\|y_i^{[k]}-Y_{\rho_i}\|_2$.  The global bound
$\Lip(M_\rho)\le CL$, together with the absolute row sums, gives
\[
 e_{k+1}\le C_ml e_k+C_mB_dl\sqrt d\,h_0^m.
\]
On each exact transported marginal, conditional expectation and the
primitive moment bound give $\|M_\rho\|_2\le CD$.  Integrating the
velocity shows $e_0\le CrD$.  Since $C_ml<1/2$ for large $d$, summing
the geometric recurrence proves \eqref{eq:probability-picard}.
The norm here is the maximum of the nodewise RMS errors; it is not the
RMS of their maximum.

For deterministic inputs the same recurrence, now in maximum node norm,
gives a bounded Lipschitz constant for every iterate.  In the correction
$y_i^{[k+1]}-y_{\mathrm{in}}$, the factor $r$ and the $CL$-Lipschitz mean
give the stronger constant $C_ml$.
\end{proof}

\subsection{The kernel action through harmonic trajectories}

The position velocity in \eqref{eq:centering-flow} contains $K(y)^TG$.
Approximating $K$ and then multiplying the error by an arbitrary momentum
would not give the global Lipschitz property needed for Picard iteration.
Instead we approximate the vector action directly, using short harmonic
trajectories and forward transport corrections.

For the kernel integral use the mesh on $[0,b_*]$ above.  Its cell lengths
also satisfy $\ell_c\ge c'h_0(1-\rho)$ for all sufficiently large $d$;
this follows from the ceiling construction because the transformed
interval length tends to infinity.  Set
$C_\rho(y)=B_\rho(y)-I$.  We first describe a circuit for
$C_\rho(y)G$ at one probability node $\rho$.

\paragraph{Backward position.}
Apply Proposition~\ref{prop:probability-circuit} backward from input
position $y$ at time $b_*$ to obtain an approximation $\widehat Y_\rho$
to $Y_\rho(y)$.

\paragraph{Short harmonic trajectories.}
At fixed probability time $\rho$, the Hamiltonian equations with
potential $H_\rho$ have position integral equation
\begin{equation}\label{eq:harmonic-integral}
 W_z=\cos(z)Y_\rho+\sin(z)G
      -r\rho\int_0^z\sin(z-v)M_\rho(W_v)\,dv.
\end{equation}
For every requested angle $z$, interpolate on the single cell $[0,z]$.
Let $z_i$ be its nodes and $\ell_j(v)$ its Lagrange basis, and set
\[
 w_{ij}^{\sin}=\int_0^{z_i}\sin(z_i-v)\ell_j(v)\,dv.
\]
Use the simultaneous Picard updates
\begin{equation}\label{eq:harmonic-picard}
 \begin{split}
 W_i^{[0]}&=\cos(z_i)\widehat Y_\rho+\sin(z_i)G,\\
 W_i^{[k+1]}&=\cos(z_i)\widehat Y_\rho+\sin(z_i)G
      -r\rho\sum_jw_{ij}^{\sin}M_\rho(W_j^{[k]}).
 \end{split}
\end{equation}
At angle zero there is no integral.  The free sine and cosine terms
are retained exactly.

\paragraph{Differentiate the transport correction.}
Let $d_j$ be the weights differentiating at zero the polynomial
interpolant on angle nodes $0,\psi,\ldots,(m-1)\psi$.  Their absolute
sum is $C_m\psi^{-1}$.  At every such angle node, compute
$T_\rho(W_z)-W_z$ by the forward probability-flow circuit, keeping only
its integral correction at the output.  For exact paths,
\begin{equation}\label{eq:kernel-action-derivative}
 \left.\frac{d}{dz}(T_\rho-\mathrm{id})(W_z)\right|_{z=0}
       =(B_\rho-I)G=C_\rho G,
\end{equation}
since $W_0=Y_\rho$ and $W'_0=G$.  The weighted sum with coefficients
$d_j$ therefore approximates $C_\rho G$ using exact mean nodes only.

Here and below the interpolation orders denoted by $m$ may be increased
independently.  Using a common sufficiently large order merely simplifies
the error displays.  For the next estimates use the same number
$N_p\ge1$ of Picard layers in the backward, harmonic, and forward
computations.

\begin{lemma}[Errors in the angle computation]
\label{lem:kernel-angle-errors}
Let $y\sim\xi_*$ and $G\sim N(0,I_d)$ independently.  The exact-input
angle differentiation in \eqref{eq:kernel-action-derivative} has RMS
truncation error at most
\begin{equation}\label{eq:angle-stencil-error}
 C_m l\sqrt d\,\psi^{m-1}R^{-2m}
       =C_ml\sqrt d\,\delta^{2tm-4t}.
\end{equation}
Before applying the angle differentiation weights, the RMS errors of the
computed forward correction values are at most
\begin{equation}\label{eq:kernel-value-error}
 C_m\polylog(d)\,l\sqrt d\,(h_0^m+\delta^{2tm})
       +(C_ml)^{N_p}CrD.
\end{equation}
Every such approximate correction is globally Lipschitz in $(y,G)$
with constant at most $C_ml$.
\end{lemma}

\begin{proof}
For exact inputs $(Y_\rho(y),G)$, harmonic motion
$\dot W=P$, $\dot P=-\nabla H_\rho(W)$ preserves
$\xi_\rho\otimes N(0,I_d)$.  The vector field is globally Lipschitz,
and its weighted divergence vanishes.  Its Lie derivative is
$\nabla^*A\nabla$ with constant skew coefficient whose $(W,P)$ block
is $-I$.  Lemma~\ref{lem:stationary-derivatives} therefore applies with
$\alpha=1$, with $a_0=l$ for $T_\rho-\mathrm{id}$, and with $a_0=L$
for $M_\rho$.  Indeed Lemma~\ref{lem:transport-derivatives} and
\eqref{eq:jets} give the required spatial bounds; changing a one-slot
split bound to a Frobenius bound costs only $\sqrt d$.
To apply the angle stencil, subtract the Taylor polynomial of
$z\mapsto(T_\rho-\mathrm{id})(W_z)$ of degree $m-1$ at zero.
Its derivative at zero is exact, and its remainder at each stencil
node has RMS norm at most $C_ml\sqrt d\,\psi^mR^{-2m}$.
Multiplication by the absolute stencil sum $C_m/\psi$ proves
\eqref{eq:angle-stencil-error}.  This uses pointwise RMS
remainders on the stationary harmonic path.

For the harmonic state equation, interpolate $M_\rho(W_v)$ in
\eqref{eq:harmonic-integral}.  Since $z\le(m-1)\psi$, its Taylor
remainder, multiplied by the prefactor $r$, is bounded by
$C_ml\sqrt d\,\psi^mR^{-2m}
=C_ml\sqrt d\,\delta^{2tm}$; the additional integration factors are
bounded and can be discarded.  The initial free-path error is at most
$CrD$, and Picard contracts by $C_ml$.  The computed harmonic state is
globally Lipschitz with bounded constant in its two inputs, by the same
row-sum argument as before.

Next compare the forward probability correction started at the exact
harmonic state.  That state has marginal law $\xi_\rho$, so
Proposition~\ref{prop:probability-circuit} applies.  Errors in its
starting state then propagate through the approximate correction with
factor $C_ml$, not a factor of order one.  This also controls the
backward-position error transmitted through the harmonic computation.
Combining these estimates proves \eqref{eq:kernel-value-error} and
the asserted global Lipschitz bound.
\end{proof}

\subsection{The zero endpoint of the kernel integral}

It remains to integrate the derivative of $C_\rho G$ against $1/\rho$.
The exact identity $B'_0=0$ removes the apparent singularity, and we
impose that identity on the interpolant as well.

On all but the first probability cell use ordinary degree-$(m-1)$
interpolation.  On the first cell use $m-1$ equally spaced value nodes
including both endpoints and the additional Hermite condition that the derivative
at zero is zero, for a total of $m$ conditions, with $m\ge3$.
This uniquely determines a polynomial of degree at most $m-1$: a
polynomial in the null space would have a double zero at zero and the
remaining $m-2$ distinct zeros.  Let $q_i$ be the resulting weights on
values when the derivative of this piecewise interpolant is integrated
against $1/\rho$.

\begin{lemma}[Derivative quadrature at zero]
\label{lem:kernel-endpoint-quadrature}
Let $y\sim\xi_*$ and $G\sim N(0,I_d)$ independently.
For exact values of $C_\rho(y)G$, the RMS error of this quadrature for
$\int_0^{b_*}B'_\rho(y)G\,d\rho/\rho$ is at most
\[
 C_m\polylog(d)\,l\sqrt d\,h_0^{m-2}.
\]
The value weights satisfy
\begin{equation}\label{eq:kernel-weight-sum}
 \sum_i|q_i|
 \le C_mh_0^{-1}(1+|\log R|+|\log h_0|).
\end{equation}
These weights are well defined even if the supplied values are
approximate: the derivative condition at zero remains exact.
\end{lemma}

\begin{proof}
By Lemma~\ref{lem:transport-derivatives} and independence of $G$,
the RMS norm of the order-$m$ derivative of $C_\rho(y)G$ is at most
$\polylog(d)l\sqrt d(1-\rho^2)^{-m}$.  Here multiplication by the
independent Gaussian converts the matrix norm to a Frobenius norm and
costs at most $\sqrt d$.

On an ordinary cell, Taylor expansion and scaling the interpolation
basis bound the derivative error by this derivative bound times
$C_m\ell_c^{m-1}$.  It is therefore at most
$C_m\polylog(d)l\sqrt d h_0^{m-1}/(1-\rho)$.
After integration against $1/\rho$, all ordinary cells contribute at
most this prefactor times
$C(1+|\log R|+|\log h_0|)$, which fits the claimed more conservative
$h_0^{m-2}$ bound.

For the first cell, subtract the Taylor polynomial at zero of degree
$m-1$.  Its derivative at zero agrees with the exact zero derivative.
The second derivative of the remainder, and that of the Hermite
interpolant of the remainder, are bounded in RMS by
$C_m\polylog(d)l\sqrt d\,\ell_c^{m-2}$.  Both first derivatives
vanish at zero.  If $E$ is the difference between the exact function
and its Hermite interpolant on this cell, then
\[
 \|E'(\rho)\|_2
 \le\int_0^\rho\|E''(v)\|_2\,dv
 \le C_mB_dl\sqrt d\,\rho\ell_c^{m-2}.
\]
Consequently the integral of $E'(\rho)/\rho$ over the first cell
has RMS norm at most $C_mB_dl\sqrt d\,\ell_c^{m-1}$.
This is within the stated, more conservative $h_0^{m-2}$ bound.

On an ordinary cell the absolute coefficient sum of the interpolant
derivative is at most $C_m/\ell_c$.  Integrating and using
$\ell_c\ge c'h_0(1-\rho)$ gives a total bounded by
\[
 C_mh_0^{-1}\int_{\rho_{\mathrm{first}}}^{b_*}
                   \frac{d\rho}{\rho(1-\rho)}.
\]
This gives the logarithms in \eqref{eq:kernel-weight-sum}.
On the first cell, scale to the unit interval.  The polynomial
derivative vanishes at zero and is bounded by distance to zero times
the second derivative, so the integral of its absolute value divided
by that distance is bounded.  Scaling back gives a weight sum
$C_m/\ell_c\le C_mh_0^{-1}$.  This proves the claim.
\end{proof}

\subsection{A globally Lipschitz approximate centering velocity}

Let $\mathrm{corr}_{ij}(y,G)$ denote the forward correction computed at
probability node $i$ and angle node $j$ by the three operations above.
The approximate position velocity is the explicit expression
\begin{equation}\label{eq:kernel-approximate-velocity}
 \widehat v_Y(y,G)=-\frac1{rs}\sum_iq_i\sum_jd_j
                         \mathrm{corr}_{ij}(y,G).
\end{equation}
Each correction retains only the integral part of the forward Picard
formula. Writing its weights as $w_{ijk}$ and its mean values as
$M_{ijk}$ gives
\[
 \mathrm{corr}_{ij}=-r\sum_kw_{ijk}M_{ijk},
\]
and hence
\begin{equation}\label{eq:kernel-scalar-expansion}
 \widehat v_Y=\frac1s\sum_{i,j,k}q_i d_j w_{ijk}M_{ijk}.
\end{equation}
The apparent factor $1/r$ has canceled algebraically. There is no
additive forward state left in this formula, and all coefficients
multiplying mean values are deterministic scalars.
Use the exact mean expression
$\widehat v_G(y,G)=(M_*(y)-M_0)/s$ for the momentum velocity.
In particular, the momentum occurs inside angle states; it is never
multiplied by a separately approximated matrix at the output.

\begin{proposition}[Kernel-action circuit]\label{prop:kernel-circuit}
On the stationary product input law, the RMS error of
$\widehat v=(\widehat v_Y,\widehat v_G)$ relative to
\eqref{eq:centering-flow} is bounded by $\polylog(d)$ times
\begin{equation}\label{eq:kernel-velocity-error}
 \begin{split}
 &\frac L\sigma\sqrt d
 \left[h_0^{m-2}+h_0^{-1}\delta^{-4t}
                      (h_0^m+\delta^{2tm})\right]\\
 &\hspace{12mm}+\delta^{-4t-w}(C_ml)^{N_p-1}
                           r\frac L\sigma D.
 \end{split}
\end{equation}
It is globally Lipschitz, with constant at most
\begin{equation}\label{eq:velocity-lipschitz}
 C_m\polylog(d)\,(L/\sigma)\psi^{-1}h_0^{-1}.
\end{equation}
For every fixed $J\ge1$, its RMS error can be made at most
$\polylog(d)\delta^J D$ with $N_p=O(J+1)$, independently of $w$
and the interpolation orders.
\end{proposition}

\begin{proof}
Combine Lemmas~\ref{lem:kernel-angle-errors} and
\ref{lem:kernel-endpoint-quadrature}, and divide by $rs$.  The value
errors are amplified only by $C_m\psi^{-1}$ and the weight sum
\eqref{eq:kernel-weight-sum}.  The exact stencil truncation has factor
$\psi^{m-1}R^{-2m}=\delta^{2tm-4t}$ and the same probability-quadrature
weight sum.  Crucially,
\[
 \frac{l}{r\sigma}=\frac L\sigma,\qquad
 \frac{l^{N_p}rD}{r\sigma}
       =l^{N_p-1}r\frac L\sigma D.
\]
Thus neither division introduces a negative power of the potentially
small scale $r$.  The resulting estimate is
\eqref{eq:kernel-velocity-error}.

Each forward correction has Lipschitz constant $C_ml$ in its inputs,
and the backward and harmonic maps have bounded constants.  Applying
the same absolute weight sums and using $l/(r\sigma)=L/\sigma$
gives \eqref{eq:velocity-lipschitz}; the momentum component has the
smaller Lipschitz bound $CL/\sigma$.

For the last assertion, increase interpolation orders until the first
line of \eqref{eq:kernel-velocity-error} is of the desired order.
Under \eqref{eq:numerical-size}, sufficiently many $O(J+1)$ Picard
layers make the second line of that order as well.  The explicit
common choice given below verifies independence from interpolation
orders and from $w$.
\end{proof}

\subsection{The mean formula and the order of parameter choices}

We can now discretize the centering flow itself.  Interpolate on the
single time cell $[0,1]$, using ordinary degree-$(m-1)$ interpolation
and partial-integral rows $w_{ij}$.  With anchor $Z_0=(Y_0,G_0)$, set
\[
 Z_i^{[0]}=Z_0,\qquad
 Z_i^{[k+1]}=Z_0+\sum_jw_{ij}\widehat v(Z_j^{[k]}).
\]
At the end keep the position components $y_j$ and output
\begin{equation}\label{eq:mean-output}
 sG_0+\sum_j\omega_jM_*(y_j),
 \qquad \sum_j|\omega_j|\le C_m,
\end{equation}
where $\omega_j$ are the full-interval integration weights.

\begin{proposition}[Mean-flow circuit]\label{prop:mean-circuit}
For ideal independent anchors $Y_0\sim\xi_*$ and $G_0\sim N(0,I_d)$,
and every fixed $J\ge1$, the circuit \eqref{eq:mean-output} can be
chosen to differ in RMS by at most
$\polylog(d)\sigma\delta^J D$ from the exact path integral
\eqref{eq:centering-output}.  The approximate joint states have RMS
errors at most $\polylog(d)\delta^JD$.  Every integration uses
$O(J+1)$ Picard layers, with a bound independent of $w$ and all
interpolation orders.
\end{proposition}

\begin{proof}
First compare the exact flow to polynomial integration of its exact
velocity.  Lemma~\ref{lem:stationary-derivatives}, applied to the state
with one additional derivative, bounds this quadrature error by
\[
 C_m\sqrt d\,\bigl((L/s)R^{-2}\bigr)^{m+1}.
\]
The approximation error of $\widehat v$ at every exact trajectory
node is bounded by Proposition~\ref{prop:kernel-circuit}, by
stationarity.  Previous state errors propagate with its global
Lipschitz constant and the absolute row sums.  The true centering
velocity need not be globally Lipschitz in the joint variables; it is
the approximate velocity's bound that is used in this recurrence.

By \eqref{eq:numerical-parameters}, $b\ge1/2$ and $w<t$ give
$b-4t-w>1/3$.  For sufficiently large $d$, fixed constants and
polylogarithms in \eqref{eq:velocity-lipschitz} can therefore be
absorbed so that the recurrence contracts by at most $\delta^{1/3}$.
Its initial error is at most $C(L/s)\sqrt d$ by
\eqref{eq:centering-initial-displacement}.  Choose the quadrature and
velocity error orders sufficiently high, and then $O(J+1)$ Picard
layers yield the state error claimed in the statement.

For the output, apply Lemma~\ref{lem:stationary-derivatives} to $M_*$
with $a_0=L$.  Its path quadrature error is at most
$C_mL\sqrt d\,((L/s)R^{-2})^m$, which can be made at most
$\polylog(d)\sigma\delta^JD$.  The additional error from evaluating
on the approximate positions is bounded by their RMS error times
$C_mL$, using the $CL$-Lipschitz bound for $M_*$.  Since $L/\sigma$
is small, this is within the same budget.
\end{proof}

We first collect the accuracy and depth estimates.  We will then expand
the actual formulas to obtain the coefficient and scale bounds required
by the recursive implementation.

\begin{corollary}[Accuracy without increasing dependency depth]
\label{prop:exact-mean-circuits}
Fix $J\ge1$.  In all the probability, harmonic, and centering
integrations above, one may take at most
\[
 N=\lceil10(J+2)\rceil
\]
Picard layers.  After any fixed $w\in(0,t)$ is chosen, a common fixed
interpolation order satisfying
\begin{equation}\label{eq:interpolation-order-choice}
 w(m-3)>J+10,\qquad t(m-3)>J+10,\qquad m>4(J+10)
\end{equation}
suffices for the following conclusions, for all sufficiently large $d$.
Probability transport from its exact starting marginal has endpoint RMS
error at most $\polylog(d)\delta^JD$; the same holds after division by
an endpoint parameter bounded away from zero.  The mean circuit with
ideal anchors has RMS error at most $\polylog(d)\sigma\delta^JD$
relative to \eqref{eq:centering-output}.

The number of layers of exact mean dependencies within either circuit
is $O((N+1)^2)$, independently of $m$, $w$, and the number of mesh cells.
All mesh nodes, interpolation weights, scalar parameters, and iteration
counts are deterministic functions of the prescribed parameters and
dimension.
\end{corollary}

\begin{proof}
For large $d$, probability and harmonic Picard contractions
$C_ml\le\polylog(d)\delta^b$ can be bounded by $\delta^{1/3}$.
The only additional inverse powers in kernel evaluation are
$\delta^{-4t-w}$, besides the divisions already canceled in
Proposition~\ref{prop:kernel-circuit}.  The centering iteration also
contracts by $\delta^{1/3}$.  The probability iteration remainder is
at most $C\delta^{N/3}D$; after the kernel weights the iteration
remainder is at most $C\delta^{N/3-4t-w}D$.  Both exponents exceed
$J$ with a fixed positive margin.  For the truncation terms it suffices
to check the exponents
\[
 wm,\quad w(m-2),\quad w(m-1)-4t,\quad
 2tm-4t-w,\quad (m+1)/3.
\]
The first four come respectively from probability quadrature, the
zero-endpoint rule, its amplified probability-value error, and the
harmonic and angle errors.  For the last one use
$(L/s)R^{-2}\le\delta^{1/3}$ for large $d$ in both outer
quadratures, dividing the mean-output error by $\sigma$ and using
$L/\sigma\le\delta^{1/3}$.
Every displayed exponent exceeds $J$ by
\eqref{eq:interpolation-order-choice}.  The outer Picard remainder
contracts by the same $\delta^{1/3}$ factor.  Thus fixed constants
and polylogarithms, including those depending on $m$, fit within the
strict exponent margins as $d\to\infty$.

Within a probability or harmonic Picard layer, all nodes depend on
only the preceding layer, regardless of the number of nodes.  A path
through one kernel evaluation passes through a backward probability
computation, a harmonic computation, and a forward probability
computation, each having at most $N$ layers.  Such a path traverses at
most $N$ layers of the outer centering iteration.  This gives depth
$O((N+1)^2)$; final quadratures add a bounded number of layers.
All weights integrate known polynomials, or polynomials against sine
or $1/\rho$ with the stated zero-endpoint condition.  They require no
oracle information and are allowable real constants in the model.
\end{proof}

The order of choices is essential: first fix the target error and the
Picard depth, then the mesh exponent $w$, and finally the interpolation
order.  The count of evaluations can depend on that order; the number
of nested dependencies will not.

\subsection{The exact circuit interface and its derivation}

To display a primitive's parameters, for $x'\in\R^d$ and
$0<r'\le1$ write
\[
 u_{x',r'}:=\int_{\R^d}g(x'+r'z)\,\nu_{x',r'}(dz).
\]
Thus $u_{x,r}=u$ for the primitive fixed above.

We specify the two formulas to be implemented.  In this subsection
$\mathcal C_M$ and $\mathcal C_S$ denote deterministic maps of their
listed anchors, with exact primitive means available as analytical
nodes.  The parameter $\sigma$ of $\mathcal C_M$ is the parent
noise parameter, so its construction sets $s=\sigma/2$ exactly once.
Define
\begin{align}
 \mathcal C_M(x,r,\sigma;y,G)
   &=\frac\sigma2G+\sum_j\omega_jM_*(y_j),
       \label{eq:exact-mean-circuit}\\
 \mathcal C_S(x,r,\eta;y)
   &=\rho_e^{-1}\widehat T_{0\to\rho_e}(y),
 &\rho_e&=(1+\eta^2/4)^{-1/2},\quad R\le\eta\le1.
       \label{eq:exact-sample-circuit}
\end{align}
The positions $y_j$ in the first line are the last centering Picard
iterates with anchors $(y,G)$, using \eqref{eq:kernel-approximate-velocity}.
The second line uses the probability formula
\eqref{eq:probability-update}.  No additional Gaussian is included in
either map.  The next section allocates additional noise and constructs
the position anchor needed for $\mathcal C_M$.

A \emph{mean dependency path} follows a mean node to a mean value used
in its center expression, and then repeats.  Direct expression size
counts occurrences in one such expression, before expanding the center
expressions of its mean nodes.

\begin{theorem}[Numerical circuit interface]
\label{thm:numerical-interface}\label{lem:compiler-circuit}
Fix $J\ge1$, the parameters \eqref{eq:numerical-parameters}, and
$N=\lceil10(J+2)\rceil$.  Choose the fixed interpolation order by
\eqref{eq:interpolation-order-choice}.  For all sufficiently large $d$,
the maps \eqref{eq:exact-mean-circuit}--\eqref{eq:exact-sample-circuit}
have the following properties.

\begin{enumerate}[label=(\roman*),leftmargin=*]
\item \emph{Deterministic scalar formulas.}
Their exact mean nodes can be ordered as $M_1,\ldots,M_Q$, where
\begin{equation}\label{eq:exact-affine-centers}
 M_i=u_{E_i,r_i},\qquad
 E_i=b_i+\sum_{j<i}a_{ij}M_j,\qquad
 E_f=b_f+\sum_ja_{fj}M_j.
\end{equation}
All $a_{ij}$ and $a_{fj}$ are deterministic scalars.  The whole node
list, its scales, and these coefficients can be formed without
knowing a mean value or making an oracle query.

\item \emph{Direct terms and scales.}
Each node has $r_i=r\tau_i$ and
\[
 \tau_i=\sqrt{1-\rho_i^2}\ge R/\sqrt5,\qquad
 b_i=x+r(\ell_i y+m_iG).
\]
For $\mathcal C_S$, the $G$ term is absent,
$\rho_i\le\rho_e$, and $(\ell_i,m_i)=(\rho_i,0)$.
For $\mathcal C_M$, $\rho_i\le b_*$ and the pair is either
$(\rho_i,0)$ or $(\rho_i\cos z,\rho_i\sin z)$, with
$|z|\le C_m\psi$.  In particular both coefficients have absolute
value at most one.  A node for $M_0=u$ has center $x$ and scale $r$.
Every mean occurring directly in the final expression for
$\mathcal C_M$ is $M_*$ at a computed position and has scale exactly
$rR$.  The final direct terms are
$b_f=\sigma G/2$ and $b_f=y/\rho_e$, respectively.

\item \emph{Coefficient sums and direct expression size.}
There are positive declared bounds $A,A_f$ such that
$\sum_{j<i}|a_{ij}|\le A$ and $\sum_j|a_{fj}|\le A_f$, namely
\begin{equation}\label{eq:coefficient-bounds}
 \begin{aligned}
 A_M&=C_1(1+\log d)^{C_1}
       \bigl[(r/\sigma)\delta^{-4t-c_1w}+r^2\bigr],
 & A_S&=C_1r^2,\\
 A_{f,M}&=C_1,
 & A_{f,S}&=C_1r,
 &c_1&=10.
 \end{aligned}
\end{equation}
Here $C_1\ge1$ depends only on the fixed parameters and orders.
Every direct expression has at most $B_d\delta^{-c_1w}$ terms.
Every mean dependency path has length at most
$D_0=3N^2+1$; using a larger absolute multiple of $(N+1)^2$ is harmless.
This depth bound is independent of $w$, $m$, and the number of mesh cells.

\item \emph{Errors on the ideal anchor laws.}
Under \eqref{eq:numerical-size}, for independent
$Y\sim\xi_*$ and $G\sim N(0,I_d)$,
\begin{equation}\label{eq:ideal-mean-circuit-error}
 W_2\bigl(\mathcal L(\mathcal C_M(x,r,\sigma;Y,G)),
                  N(u,\sigma^2I_d/4)\bigr)
       \le B_d\sigma\delta^JD(x).
\end{equation}
For $\mathcal C_S$ only $l\le B_d\delta^b$, $b\ge1/2$, is needed.
With $Y\sim N(0,I_d)$,
\begin{equation}\label{eq:ideal-sample-circuit-error}
 W_2\bigl(\mathcal L(\mathcal C_S(x,r,\eta;Y)),
                    \nu_{x,r}*N(0,\eta^2I_d/4)\bigr)
       \le B_d\delta^JD(x).
\end{equation}
The coefficient and dependency assertions are deterministic, for every
choice of anchors.  The distributional estimates require the specified
ideal anchor laws.
\end{enumerate}
\end{theorem}

\begin{proof}
We expand only the additive state formulas, stopping whenever an earlier
mean value is reached.  This convention is what makes the coefficient
bounds local: the center expression inside that earlier mean is kept as
its own node.

\emph{The three transport formulas.}
A probability Picard state has its input position as direct term and
mean coefficients with absolute sum at most $C_mr$.  A harmonic state
has direct term $\cos z$ times its input position plus $\sin z$ times
its input momentum, with mean coefficients of absolute sum at most
$C_mr$.  The sine-weighted row sums are bounded because their integration
intervals have bounded length.  These assertions hold at every iteration:
its additive direct term is the original input, not the preceding
Picard iterate.

For a kernel evaluation, the scalar expansion
\eqref{eq:kernel-scalar-expansion} has absolute coefficient sum at most
$B_d\sigma^{-1}\delta^{-4t-w}$, by the derivative-quadrature,
angle-stencil, and forward-integral weight bounds.  The momentum
formula $(M_*-M_0)/s$ has coefficient sum at most $C/\sigma$.
A centering Picard update has the original anchor as direct term and
bounded outer integral rows.  Its position and momentum corrections
thus have the same respective coefficient bounds.

\emph{Centers and their coefficients.}
Consider a mean evaluation inside any of these formulas.  Its physical
center is $x+r\rho v$ when its normalized position is $v$, by
\eqref{eq:conditional-tilt}.  Substitute its direct starting states:
a backward probability state starts at the current centering position;
a harmonic state starts at that backward position and the current
momentum; and a forward probability state starts at the harmonic
position.  Stop each substitution at the mean terms just described.
Only the free terms pass through this chain.  Thus their anchor
coefficients are exactly $(\rho,0)$ or
$(\rho\cos z,\rho\sin z)$.  For sampling there is only the first
probability computation, so the coefficient is $(\rho,0)$.
Multiplication by $r\rho$, with $\rho\le1$, bounds the center
coefficient sum in the mean case by
\[
 B_d(r/\sigma)\delta^{-4t-w}+C_mr^2.
\]
The term $C_mr^2$ accounts for the backward, harmonic, and forward
integral corrections.  In the sampling case the bound is simply
$C_mr^2$.  Enlarging these positive upper bounds to
\eqref{eq:coefficient-bounds} is convenient for the implementation.
The mean output is the single quadrature in
\eqref{eq:exact-mean-circuit}, so its coefficient sum is at most
$C_m$ and its means all have probability time $b_*$.  For sampling,
dividing a probability correction by $\rho_e\ge2/\sqrt5$ gives
output coefficient sum at most $C_mr$.  These observations also give
the stated final direct terms.

The scale at probability time $\rho$ is
$r\sqrt{1-\rho^2}$.  In the mean computation $\rho\le b_*$, giving
$\tau\ge R$, including equality for the output means.  In the
sampling computation,
\[
 \sqrt{1-\rho_e^2}=\frac\eta{\sqrt{4+\eta^2}}
                 \ge\frac R{\sqrt5}.
\]
This proves every assertion about the anchors and scales.

\emph{Expression size and depth.}
Let $P\le B_d\delta^{-w}$ bound the number of nodes in any one
probability mesh.  Angle stencils, harmonic quadratures, and outer
centering quadratures have fixed numbers of nodes once $m$ is fixed.
Formula \eqref{eq:kernel-scalar-expansion} has at most $C_mP^2$
mean terms: one factor $P$ comes from probability-time derivative
quadrature, and the other from the forward transport integral.
A direct centering position has the same bound, and its momentum has
only a fixed number of terms.  Substituting their free states into a
backward, harmonic, or forward center adds at most one backward list,
one harmonic list, and one forward list.  No substitution continues
through a mean value.  Hence $B_d\delta^{-2w}$ already bounds the
direct expression size; the exponent $c_1w=10w$ is a uniform
conservative choice.

All Picard updates at a fixed layer depend only on earlier layers.
Order an internal kernel computation as backward transport, harmonic
motion, and forward transport, after its input states.  A mean path
within this computation adds at most $3N$ nodes above a path in those
inputs.  There are at most $N$ outer Picard layers, and one final
mean layer in \eqref{eq:exact-mean-circuit}.  Thus $3N^2+1$ bounds
the depth.  The sampling path has depth at most $N$.  This gives an
acyclic order for \eqref{eq:exact-affine-centers}.  Meshes, weights,
angles and partial integral rows depend only on the prescribed scalar
parameters; no exact mean is needed to form any of them.

\emph{Ideal-law errors.}
Couple the numerical mean formula and the exact centering path integral
by the same ideal anchors.  Corollary~\ref{prop:exact-mean-circuits}
bounds their RMS difference by $B_d\sigma\delta^JD(x)$, while
Proposition~\ref{prop:centering-law} identifies the exact output as
$N(u,\sigma^2I_d/4)$.  This coupling proves
\eqref{eq:ideal-mean-circuit-error}.  For sampling, the exact
probability transport from a standard Gaussian has law $\xi_{\rho_e}$.
After division by $\rho_e$ this is the law of
\[
 Z+\frac{\sqrt{1-\rho_e^2}}{\rho_e}G
       =Z+\frac\eta2G,
 \qquad Z\sim\nu_{x,r},\quad G\sim N(0,I_d)\text{ independent}.
\]
The same-anchor RMS estimate and $\rho_e\ge2/\sqrt5$ prove
\eqref{eq:ideal-sample-circuit-error}.
\end{proof}

For the implementation we use $J=K_0+2$.  The theorem fixes $N$ and
$D_0$ before $w$; the implementation can therefore choose $w$ using
its recursion count, and then choose $m$ by
\eqref{eq:interpolation-order-choice}.  The formulas above introduce
no random failure event.  Their errors are RMS estimates under the
ideal input laws, and their algebraic structure holds for every input.
